\phantomsection\label{fig:vol04:liao-western-xia:multi-polity-frontiers}
\begin{center}
\begin{tikzpicture}[x=.88cm,y=.88cm,font=\small]
  \fill[paper] (0,0) rectangle (11.7,6.35);

  \fill[source!20] (.55,3.85) -- (11.15,3.6) -- (11.1,5.85) -- (.75,5.85) -- cycle;
  \node[align=center] at (8.8,4.9) {辽\\燕云、草原与东北\\的多重政治中心};
  \fill[timeline!22] (.55,1.0) -- (3.35,.8) -- (3.65,3.35) -- (.65,3.55) -- cycle;
  \node[align=center] at (1.95,2.15) {西夏\\河西、鄂尔多斯与\\边城网络};
  \fill[dynasty!19] (3.9,.55) -- (11.2,.75) -- (11.1,3.25) -- (3.75,3.45) -- cycle;
  \node[align=center] at (8.1,1.65) {北宋\\河北、陕西与\\中原财政腹地};

  \draw[very thick,color=timeline!85] (.75,3.62) .. controls (3.05,3.9) and (5.35,3.12) ..
    (7.75,3.55) .. controls (9.25,3.82) and (10.1,3.48) .. (11.05,3.62);
  \node[draw=timeline!80,fill=paper,rounded corners=2pt,align=center,inner sep=3pt]
    at (6.05,3.67) {盟约、堡寨、边市与\\使节往来的接触地带};

  \fill[dynasty] (7.55,2.1) circle (.12);
  \node[below,align=center] at (7.55,2.02) {东京开封\\宋廷中枢};
  \fill[timeline] (5.65,3.08) circle (.12);
  \node[below,align=center] at (5.65,2.98) {澶州\\1004年对峙};
  \fill[source] (8.85,4.5) circle (.12);
  \node[above,align=center] at (8.85,4.58) {燕云城市\\与道路节点};
  \fill[timeline!85] (2.2,3.0) circle (.12);
  \node[left,align=center] at (2.1,3.0) {灵州、兴庆\\方向};

  \draw[<->,thick,dashed,color=timeline!90] (5.15,3.13) -- (6.35,3.14);
  \node[text=timeline!90!black,above,align=center,font=\scriptsize] at (5.75,3.12)
    {谈判与互市};
  \draw[->,ultra thick,color=source!85] (8.55,4.38) ..
    controls (7.6,3.85) and (6.85,3.55) .. (5.9,3.22);
  \node[text=source,align=center,font=\scriptsize] at (7.55,4.12)
    {辽军南下\\方向（1004）};
  \draw[->,thick,color=timeline!85] (2.42,2.88) ..
    controls (3.5,2.8) and (4.05,2.48) .. (4.85,2.4);
  \node[text=timeline!85!black,align=center,font=\scriptsize] at (3.55,2.15)
    {宋夏战事、和议\\与西北调度};

  \node[draw=source!75,fill=source!12,rounded corners=2pt,align=center,inner sep=3pt]
    at (1.95,4.92) {女真／金的兴起\\1114--1125年改变\\东北与燕云关系};
  \draw[->,thick,dashed,color=source!80] (3.2,4.75) -- (6.9,4.65);
  \node[text=source,align=center,font=\scriptsize] at (8.95,.76)
    {疆界、控制范围与路线均为约略示意；不同年份的关系不可叠作同一时点。};
\end{tikzpicture}

\smallskip
\small\textit{图\quad 十一至十二世纪宋、辽、西夏与金兴起前后的多政权边境关系示意：据《宋史》〈真宗本纪〉、〈夏国传〉、〈辽史〉〈圣宗本纪〉及营卫、地理诸志，《续资治通鉴长编》景德、庆历相关条，《宋会要辑稿·蕃夷》，并参照《金史》〈太祖本纪〉与《三朝北盟会编》所见1114--1125年东北、燕云的转换。图示澶渊对峙、宋夏接触、边市和金之兴起如何使边境成为军队、文书、贸易与城镇共同构成的地带；它不提供精确国界、兵力、税额或各城持续控制的断言。}
\end{center}
